\documentclass[11pt]{scrartcl}
%\includeonly{chapters/fundamentals-of-inequalities}
\usepackage[sexy]{evan}

\title{NOTI A CHI?}
\author{Anto Turneight}
\date{\today}
\begin{document}

\maketitle
 \begin{abstract}
Questo è un mio tentativo di enucleare in un pdf alcune delle possibili
soluzioni dei famosissimi NOTI del senior. Per chi non sapesse bene cosa siano,
sono $10$ problemi olimpici, alcuni dei quali vengono proposti nel test finale
dello stage senior, che penso abbiano lo scopo di far prendere familiarità
con alcune tecniche standard utili per la risoluzione di diversi problemi di
natura olimpica. Essendo problemi standard, secondo me non sono il massimo del
divertimento, ma sono comunque utili anche per migliorare le proprie abilità
di proof-writing, che credo sia il motivo per cui li mettano nel test finale.
La cosa che mi ha sorpreso è che in rete non sono riuscito a trovare un
pdf che raccoglie le soluzioni, e quindi volevo porre rimedio a sta cosa.
Ci tengo però a precisare che su youtube si trovano i video di Giuseppe Bassi, dove
viene spiegato bene il procedimento.
Le soluzioni possono contenere errori e non essere le migliori di sempre,
sono per lo più mie integrate con qualche idea delle sol di Giuseppe Bassi,
quindi correzioni e suggerimenti sono bene accetti. Spero
possano essere utili a qualcuno. Il template latex è quello solito di
Evan Chen. I problemi sono raggruppati nelle loro quattro aree tematiche (Algebra
,Combinatoria, Geometria, teoria dei Numeri).
 \end{abstract}

 \tableofcontents


\newpage


\section{Problemi}
\begin{enumerate}[\bfseries 1.]
\item %% Problem 1
  (A1)
Un polinomio $P(z)$ a coefficienti complessi ha grado 2002 e ha le
radici (complesse) a due a due distinte.
Dimostrare che esistono numeri complessi $a_1, \ldots, a_{2002}$
tali che, se si definiscono i polinomi $P_n(z)$ mediante la ricorrenza
\[
  P_1(z)=z-a_1, \quad P_{n+1}(z)=\left[P_n(z)\right]^2-a_{n+1},
\]
allora $P(z)$ divide $P_{2002}(z)$.

\item %% Problem 2
  (A2)
10. Siano $a, b, c$ numeri reali positivi. Dimostrare che
$$
\frac{a}{b+c}+\frac{b}{c+a}+\frac{c}{a+b} \geq \frac{3}{2}
$$
in almeno tre modi diversi, usando
- la sostituzione $a+b=x, b+c=y, c+a=z$;
- Cauchy Schwarz;
- le disuguaglianze di convessità.

Determinare quindi le migliori costanti $C_1, C_2$ tali che

\[
  C_1 \leq \frac{a}{a+b}+\frac{b}{b+c}+\frac{c}{c+a} \leq C_2 .
\]

\item %% Problem 3
  (A3)
Trovare tutte le funzioni $f: \mathbb{R} \rightarrow \mathbb{R}$ tali che
$$
f(x f(x)+f(y))=[f(x)]^2+y
$$
per ogni $x, y$ reali.

\item %% Problem 4
  (C1) Ad un party prendono parte $12 k$ persone.
Ciascuna di esse stringe la mano ad esattamente $3 k+6$ persone.
Si sa che esiste un numero $N$ tale che, per ogni coppia di persone
$A, B$, il numero di invitati che stringe la mano sia ad $A$ sia a $B$
è esattamente $N$.
Determinare per quali valori (interi) di $k$ si può realizzare
questa situazione.

\item %% Problem 5
  (C2)Alberto e Barbara tracciano un grafo su di un foglio ed iniziano
a giocare al seguente gioco. A turno, iniziando da Alberto,
ciascun giocatore effettua una delle seguenti due mosse:
- cancellare tre segmenti che formano un triangolo;
- dati tre punti di cui due non collegati, ma collegati al terzo,
cancellare i due collegamenti presenti e unire i due punti non collegati.

Il giocatore che non può fare mosse valide perde la partita.
Dimostrare che l'esito della partita non dipende da come giocano
Alberto e Barbara, e determinare un criterio per stabilire chi vince
in funzione della configurazione iniziale.
\item %% Problem 6
  (G1)Sia dato un triangolo $A B C$ e si fissino i punti $A^{\prime}, B^{\prime}, C^{\prime}$
sui lati opposti ai vertici $A, B$, $C$, rispettivamente,
in modo che le rette $A A^{\prime}, B B^{\prime}, C C^{\prime}$
siano concorrenti in un punto $P$ interno al triangolo.
Sia $d$ il diametro del cerchio circoscritto al triangolo $A B C$,
e sia $S^{\prime}$ l'area del triangolo $A^{\prime} B^{\prime} C^{\prime}$.
Dimostrare che
$$
d \cdot S^{\prime}=A B^{\prime} \cdot B C^{\prime} \cdot C A^{\prime} .
$$.

\item %% Problem 7
  (G2)Siano $A B M N$ e $B C Q P$ i quadrati costruiti sui lati
$A B$ e $B C$ di un triangolo, esternamente al triangolo stesso.
Dimostrare che i centri di tali quadrati ed i punti medi di $A C$ ed
$M P$ sono i vertici di un quadrato.

\item %% Problem 8
  (G3) Sia $M N$ una retta parallela al lato $B C$ di un triangolo $A B C$ che
interseca i lati $A B$ e $A C$ in $M$ ed $N$ rispettivamente.
Sia $P$ l'intersezione delle rette $B N$ e $C M$, e si supponga
che le circonferenze circoscritte ai triangoli $B M P$ e $C N P$
abbiano un ulteriore punto d'intersezione $Q$, oltre a $P$.
Mostrare che $\angle BAQ=\angle CAP$.
\item %% Problem 9
  (N1) Per ogni intero positivo $n$, sia $d_n$ il massimo comun divisore tra
\[
100 + n^2 \qquad \text{e} \qquad 100 + (n+1)^2 .
\]

Determinare il massimo valore possibile per $d_n$.

\item %% Problem 10
  (N2) Consideriamo l’insieme

\[
D = \{\, n \in \mathbb{N} : n \mid (2^n + 1) \,\}.
\]
\begin{enumerate}[(a)]
    \item Determinare tutti i primi che appartengono a \(D\).
    \item Determinare tutte le potenze di primi che appartengono a \(D\).
    \item Determinare tutti gli elementi di \(D\) che sono prodotto di due primi.
    \item Dimostrare che tutti gli elementi di \(D\) sono multipli di \(3\).
    \item Determinare tutti gli elementi di \(D\) della forma \(p^2 q\), con \(p\) e \(q\) primi distinti.
\end{enumerate}

\end{enumerate}




\pagebreak




\section{SOLUZIONI AI PROBLEMI}
\section{A}

\subsection{A1}
\begin{lemma*}
Dati $n\geq 1$ e $n+1$ valori complessi arbitrari $x_1,\dots,x_{n+1}$, esistono $a_1,\dots,a_n\in\mathbb{C}$ tali che $P_n(x_1)^2=\dots=P_n(x_{n+1})^2$.
\end{lemma*}

\begin{proof}
Per $n=1$: si prende $a_1=\frac{x_1+x_2}{2}$. Allora $P_1(x_1)=-P_1(x_2)$, dunque i quadrati coincidono.

Sia $n\geq 2$ e si assuma il lemma vero per $n-1$. Per ipotesi induttiva esistono $a_1,\dots,a_{n-1}$ tali che
$$
P_{n-1}(x_1)^2=\dots=P_{n-1}(x_n)^2=:c.
$$
Sia $A:=P_{n-1}(x_{n+1})$. Si sceglie
$$
a_n=\frac{A^2+c}{2}.
$$
Allora
$$
P_n(x_i)=P_{n-1}(x_i)^2-a_n=c-a_n=-\bigl(A^2-a_n\bigr)=-P_n(x_{n+1})
$$
per $i=1,\dots,n$, quindi $P_n(x_1)^2=\dots=P_n(x_{n+1})^2$.
\end{proof}

Siano $\alpha_1,\dots,\alpha_{2002}$ le radici di $P$, distinte. Basta costruire $a_1,\dots,a_{2002}$ tali che $P_{2002}(\alpha_i)=0$ per ogni $i$.

Applicando il lemma con $n=2001$ e $\{x_i\}=\{\alpha_i\}$ si ottengono $a_1,\dots,a_{2001}$ tali che
$$
P_{2001}(\alpha_1)^2=\dots=P_{2001}(\alpha_{2002})^2=:c.
$$
Si pone infine $a_{2002}=c$. Allora $P_{2002}(\alpha_i)=0$ per ogni $i$, dunque $P$ divide $P_{2002}$.



\subsection{A2}
\textbf{Modo 1 (sostituzione).} Posto $x=b+c,\;y=c+a,\;z=a+b$, si ha
$$
\sum_{\text{cyc}}\frac{a}{b+c}=\frac12\sum_{\text{cyc}}\frac{y+z-x}{x}
=\frac12\Bigl(\sum_{\text{sym}}\frac{x}{y}-3\Bigr)\ge\frac12(6-3)=\frac32,
$$
poich\'e $\frac{x}{y}+\frac{y}{x}\ge2$ per ogni coppia.

\textbf{Modo 2 (Cauchy--Schwarz).} Per Titu,
$$
\sum_{\text{cyc}}\frac{a}{b+c}=\sum_{\text{cyc}}\frac{a^2}{a(b+c)}
\ge\frac{(a+b+c)^2}{2(ab+bc+ca)}\ge\frac{3(ab+bc+ca)}{2(ab+bc+ca)}=\frac32,

$$
dove $(a+b+c)^2\ge3(ab+bc+ca)\iff a^2+b^2+c^2\ge ab+bc+ca$.

\textbf{Modo 3 (convessit\`a).} Per omogeneit\`a, poniamo $a+b+c=1$; allora $b+c=1-a$ etc.
$$
\sum_{\text{cyc}}\frac{a}{1-a}\ge\frac32.
$$
$f(x)=\frac{x}{1-x}$ su $(0,1)$ ha $f''(x)=\frac{2}{(1-x)^3}>0$; per Jensen,
$$
\frac{f(a)+f(b)+f(c)}{3}\ge f\!\left(\frac13\right)=\frac12.
$$

\textbf{Determinazione delle costanti.} Sia $S=\frac{a}{a+b}+\frac{b}{b+c}+\frac{c}{c+a}$.
Per il limite inferiore, $\frac{a}{a+b}>\frac{a}{a+b+c}$ e cicliche, quindi $S>1$.
Con $a=1,\;b=t,\;c=t^2$ e $t\to\infty$ si ha $S\to1$, dunque $C_1=1$.
Per il limite superiore, $\frac{a}{a+b}<\frac{a+c}{a+b+c}$ etc., quindi $S<2$.
Con $a=\varepsilon,\;b=\varepsilon^2,\;c=1$ e $\varepsilon\to0^+$ si ha $S\to2$, dunque $C_2=2$.



\subsection{A3}

Denotiamo con $P(x,y)$ l'asserzione
\[
f\bigl(x\,f(x)+f(y)\bigr)=[f(x)]^2+y.
\]
Da $P(0,y)$ segue $f(f(y))=[f(0)]^2+y$. Dunque $f$ è biiettiva e esiste $a$ con $f(a)=0$.
Da $P(a,0)$ si ottiene $f(f(0))=0$. Ma $f(f(0))=[f(0)]^2$, quindi $f(0)=0$ e $f(f(y))=y$ per ogni $y$.
Sostituendo $x$ con $f(x)$ in $P(x,y)$ e usando $f(f(x))=x$ si ha
\[
f\bigl(x\,f(x)+f(y)\bigr)=x^2+y.
\]
Confrontando con $P(x,y)$ si deduce $[f(x)]^2=x^2$, cioè $f(x)=\pm x$ per ogni $x$.
Le funzioni $f(x)=x$ e $f(x)=-x$ funzionano. Se esistessero $a,b\neq0$ con $f(a)=a$
e $f(b)=-b$, allora $P(a,b)$ darebbe $f(a^2-b)=a^2+b$, impossibile
poiché il primo membro vale $\pm(a^2-b)$ ed $a$ e $b$ sono non nulli.

Quindi le sole soluzioni sono $f(x)=x$ e $f(x)=-x$.


\section{C}

\subsection{C1}
Il grafo $G$ delle strette di mano $v=12k$ vertici ed è regolare di grado $d=3k+6$.
Inoltre esiste un intero $N$ tale che ogni coppia di vertici distinti
ha esattamente $N$ vicini comuni.
Contiamo in due modi le terne ordinate $(A,B,C)$ di vertici distinti con $A\sim B$ e $A\sim C$.
Fissato $A$, si possono scegliere $B$ e $C$ in $\binom{d}{2}$ modi,
quindi il totale è $v\binom{d}{2}$.

D’altra parte, fissata una coppia $\{B,C\}$, il numero di $A$ comuni è $N$,
e le coppie sono $\binom{v}{2}$. Si ottiene perciò
\[
12k\cdot\binom{3k+6}{2}=N\cdot\binom{12k}{2},
\]
da cui
\[
N=\frac{(3k+6)(3k+5)}{12k-1}.
\]
Dunque $12k-1$ divide $(3k+6)(3k+5)$. Poiché $\gcd(12k-1,3)=1$, segue che
$12k-1$ divide $3k^2+11k+10$. Eseguendo la divisione euclidea di $3k^2+11k+10$
per $12k-1$ si trova un resto costante, e si conclude che $12k-1$ divide $175$.
I divisori positivi di $175$ sono $1,5,7,25,35,175$.
L’unica soluzione intera $k\ge 1$ per cui $12k-1$ è uno di questi divisori e
$N$ risulta intero positivo è $k=3$, che dà $N=6$. \\
\textbf{Costruzione per $k=3$}.
I $36$ vertici li immaginiamo come le celle di una schacchiera $6 \times 6$ .
Due punti distinti sono adiacenti se e solo se coincidono in esattamente
una delle tre coordinate
\[
(x,y)\mapsto x,\qquad (x,y)\mapsto y,\qquad (x,y)\mapsto x+y\bmod 6.
\]
Le tre partizioni indotte (righe, colonne, somme costanti modulo 6) sono a due a due ortogonali:
cioè una classe di una partizione e una classe di un’altra si intersecano in un solo punto.
Di conseguenza ogni vertice appartiene a esattamente una classe di ciascuna partizione,
e le tre classi contribuiscono $5+5+5$ punti distinti; il grado è dunque $15$.\\
Siano ora $A$ e $B$ due vertici distinti. Un terzo punto $C$ è vicino comune se e solo se coincide con
$A$ in una coordinata e con $B$ in una coordinata.\\
- Se $A$ e $B$ condividono già una coordinata, gli altri $4$ punti della retta comune e
le due combinazioni “incrociate” delle restanti coordinate (ciascuna delle quali,
per ortogonalità, determina un unico punto) danno esattamente $6$ vicini comuni. \\
- Se $A$ e $B$ non condividono nessuna coordinata, le $3\times 2=6$ possibili scelte
di coordinate determinano, sempre per ortogonalità, esattamente $6$ punti distinti.
In entrambi i casi il numero di vicini comuni è $6$.
\begin{remark*}
La parte non standard del problema è la costruzione, ed è anche la parte
più interessante secondo me. Infatti è un topos abbastanza ricorrente quello di
dare un interpretazione circa matriciale al grafo, ad esempio TF senior 25 problema
dimostrativo di C usa un'idea simile per la costruzione (anche se è un paragone un po' tirato per i capelli).
Facendo un peletto in più di ricerca (chat
gpt), si scopre che il grafo è uno strongly regular graph con parametri
\[ \operatorname{srg}(36,15,6,6) \]
(e quindi con $  \lambda=\mu  $).
Più specificamente, all’interno di questa classe, si tratta del grafo di Cayley del
gruppo abeliano  $  \mathbb{Z}_6\times\mathbb{Z}_6  $ rispetto all’insieme
di connessione costituito dagli elementi non banali dei tre sottogruppi
\[\langle(1,0)\rangle,\qquad\langle(0,1)\rangle,\qquad\langle(1,-1)\rangle.\]
Equivalentemente, è il grafo ottenuto come unione delle cricche provenienti
da tre classi parallele a due a due ortogonali (righe, colonne e somme costanti)
su un insieme di $  6^2  $ punti.
\end{remark*}

\subsection{C2}
Supponiamo il grafo abbia $n$ vertici di grado $d_1,d_2 \ldots , d_n$. Innanzitutto
notiamo che il numero di archi del grafo decresce strettamente dopo ogni mossa e
quindi il gioco deve finire. Consideriamo la quantità
\[T = \sum_{i=1}^{n} \lfloor \frac{d_i}{2} \rfloor\].
L'osservazione principale è che $T$ dopo ogni mossa decresce sempre o
 di $3$ nel caso vengano tolti tre archi di un triangolo, in queanto il grado
 di ogni vertice che lo forma decresce di $2$,oppure di $1$ quando, la seconda
 mossa viene applicata, infatti due vertici del triangolo mantengono lo stesso
 grado mentre l'altro decresce di $2$. Osservo inoltre che una mossa è possibile
 se e solo se $T>0$. Infatti se esiste un vertice di grado due o superiore posso
 applicare quantomeno la seconda mossa, se invece non esistesse nessun vertice di
 grado almeno due, entrambe le mosse sono precluse. Di conseguenza il gioco termina
 nel momento in cui $T=0$. Poichè $T$ decresce sempre di un numero dispari,
 la parità del numero di mosse possibili è unicamente determinata, e di conseguenza,
 è determinato anche l'esito del gioco, infatti se $T$ è disparti vince Alberto,
 altrimenti Barbara.

\section{G}

\subsection{G1}
Siano $a=BC$, $a_1=BA'$, $a_2=A'C$, $b=CA$, $b_1=CB'$, $b_2=B'A$, $c=AB$, $c_1=AC'$, $c_2=C'B$.

Per il teorema di Ceva (le ceviane sono concorrenti) si ha
\[
\frac{a_1}{a_2}\cdot\frac{b_1}{b_2}\cdot\frac{c_1}{c_2}=1,
\]
da cui $a_1 b_1 c_1=a_2 b_2 c_2=x$.


\begin{center}
\begin{asy}
import graph;
size(5cm);
defaultpen(linewidth(0.8));
pen tri = rgb(0,0.8,0);
pen thick = rgb(0,0.4,0.6);

pair A  = (3.5761,24.9047);
pair B  = (1,17);
pair C  = (10,17);

pair Ap = (4.4565,17);
pair Bp = (6.0339,21.8804);
pair Cp = (2.2920,20.9645);

pair X  = (4.1196,20.0245);

draw(A--B--C--cycle, tri+linewidth(0.5));
draw(Ap--Bp--Cp--cycle, tri+linewidth(0.5));
draw(Ap--A, thick+linewidth(0.6));
draw(B--Bp, thick+linewidth(0.6));
draw(C--Cp, thick+linewidth(0.6));
draw(B--Ap);
draw(C--Ap);
draw(C--Bp);
draw(A--Bp);
draw(A--Cp);
draw(B--Cp);

dot(A);  dot(B);  dot(C);
dot(Ap); dot(Bp); dot(Cp);
dot(X);

label("$A$",  A,  NE);
label("$B$",  B,  SW);
label("$C$",  C,  SE);

label("$A'$", Ap, S);
label("$B'$", Bp, NE);
label("$C'$", Cp, NW);
label("$X$",  X,  NE);

label("$a_1$", midpoint(B--Ap), NE);
label("$a_2$", midpoint(C--Ap), NE);

label("$b_1$", midpoint(C--Bp), NE);
label("$b_2$", midpoint(A--Bp), NE);

label("$c_1$", midpoint(A--Cp), NE);
label("$c_2$", midpoint(B--Cp), NE);

\end{asy}

\end{center}
Denotiamo con $[XYZ]$ l’area del triangolo $XYZ$ e sia $S=[ABC]$. L’area $S'=[A'B'C']$
si ottiene sottraendo da $S$ le aree dei tre triangoli $BA'C'$, $CB'A'$ e $AC'B'$:
\[
S'=S-[BA'C']-[CB'A']-[AC'B'].
\]
Osserviamo che
\[
[BA'C']=\frac{a_1}{a}\cdot[BCC']=\frac{a_1}{a}\cdot\frac{c_2}{c}\cdot S,
\]
e analogamente
\[
[CB'A']=\frac{b_1}{b}\cdot\frac{a_2}{a}\cdot S,\qquad
[AC'B']=\frac{c_1}{c}\cdot\frac{b_2}{b}\cdot S.
\]
Sostituendo si ottiene
\[
S'=S\Biggl(1-\frac{a_1 c_2}{a c}-\frac{b_1 a_2}{b a}-\frac{c_1 b_2}{c b}\Biggr)
=S\cdot\frac{abc-a b_2 c_1-b c_2 a_1-c a_2 b_1}{abc}.
\]
Ricordiamo la formula $S=\dfrac{abc}{2d}$ (equivalente a $S=\dfrac{abc}{4R}$ con $d=2R$).
Moltiplicando per $d$ si ha
\[
d\cdot S'=\frac{abc-a b_2 c_1-b c_2 a_1-c a_2 b_1}{2}.\tag{$*$}
\]
Vogliamo mostrare che il membro destro di $(*)$ coincide con $a_2 b_2 c_2$.
Sostituiamo $a=a_1+a_2$, $b=b_1+b_2$, $c=c_1+c_2$ e sviluppiamo il numeratore:
\begin{align*}
abc-a b_2 c_1-b c_2 a_1-c a_2 b_1
&=(a_1+a_2)(b_1+b_2)(c_1+c_2)\\
&\quad-(a_1+a_2)b_2 c_1-(b_1+b_2)c_2 a_1-(c_1+c_2)a_2 b_1.
\end{align*}
Espandendo e semplificando si ottiene esattamente $ d \cdot S'
\frac{a_1 b_1 c_1+a_2 b_2 c_2}{2}=x$.
Inoltre $x=a_2 b_2 c_2=AB'\cdot BC'\cdot CA'$, da cui la tesi segue.
\subsection{G2}
Sia $L$ il punto medio di $CA$, sia $K$ il punto medio di $MP$, e siano $O_1$ e $O_2$
i centri dei quadrati $BCQP$ e $ABMN$ rispettivamente.
Per definizione $O_1$ è il punto medio di $CP$ e $O_2$ è il punto medio di $AM$.

Consideriamo la rotazione $\theta$ di centro $A$ e angolo $90^\circ$ in
senso antiorario, quindi  $\theta(M)=A$ e $\theta(C)=P$.

Allora $\theta$ manda il segmento $MC$ nel segmento $AP$, quindi $MC=AP$ e
$PA\perp CM$.
Adesso consideriamo il quadrilatero $ACPM$, i punti $L$, $O_1$, $K$, $O_2$
sono i punti medi di $AC$, $CP$, $PM$ e $MA$ rispettivamente.
Ne segue
\[
LO_1=KO_2=\frac12 AP=\frac12 CM=O_1K=O_2L.
\]
Dunque $LO_1KO_2$ è un rombo. Inoltre, poiché $MC\perp AP$, per Talete si ha
\[
LO_1\parallel KO_2\parallel AP\perp CM\parallel O_1K\parallel O_2L.
\]
I lati consecutivi sono perciò perpendicolari, e $LO_1KO_2$ è dunque un quadrato.

\begin{center}
\begin{asy}
import graph;
size(8cm);
defaultpen(linewidth(0.8));
pen thinGreen = rgb(0,0.8,0) + linewidth(0.6);
pen medRed   = rgb(1,0,0.2) + linewidth(0.9);
pen medCyan  = rgb(0.2,1,0.8) + linewidth(0.6);
pen thickBlue = rgb(0,0,0.8) + linewidth(0.9);
pen aquaThin = rgb(0,1,1) + linewidth(0.6);
pen dotstyle = black;

pair A = (4.,19.);
pair B = (5.444732812200629,23.55713789950519);
pair C = (11.,19.);
pair M = (0.8875949126954406,25.00187071170582);
pair N = (-0.5571378995051903,20.44473281220063);
pair Q = (15.557137899505188,24.555267187799377);
pair P = (10.001870711705816,29.112405087304563);
pair L = (7.5,19.);
pair K = (5.444732812200629,27.057137899505193);
pair O1 = (10.500935355852908,24.056202543652283);
pair O2 = (2.4437974563477205,22.00093535585291);

draw(A--B--C--cycle, thinGreen);
draw(A--B--M--N--cycle, thinGreen);
draw(B--C--Q--P--cycle, thinGreen);
draw(M--B--C--cycle, medCyan);
draw(A--B--P--cycle, aquaThin);

draw(A--C, linewidth(1.2));
draw(M--P, linewidth(1.2));
draw(A--M, linewidth(1.2));
draw(C--P, linewidth(1.2));

draw(O2--L, medRed);
draw(O1--L, medRed);
draw(O2--K, medRed);
draw(O1--K, medRed);

draw(M--B, medCyan);
draw(B--C, medCyan);
draw(C--M, thickBlue);
draw(A--B, aquaThin);
draw(B--P, aquaThin);
draw(P--A, thickBlue);

dot(A, dotstyle+linewidth(2));
dot(B, dotstyle+linewidth(2));
dot(C, dotstyle+linewidth(2));
dot(M, dotstyle+linewidth(2));
dot(N, dotstyle+linewidth(2));
dot(Q, dotstyle+linewidth(2));
dot(P, dotstyle+linewidth(2));
dot(L, dotstyle+linewidth(2));
dot(K, dotstyle+linewidth(2));
dot(O1, dotstyle+linewidth(2));
dot(O2, dotstyle+linewidth(2));

label("$A$", (3.558128578939537,18.280096255568794), NE*0.5);
label("$B$", (5.121314943641582,24.096167680436928), NE*0.5);
label("$C$", (11.06861019325511,18.365223058407362), NE*0.5);
label("$M$", (0.18020861843396907,25.22813022039358), NE*0.5);
label("$N$", (-1.7,20.215153257728403), NE*0.5);
label("$Q$", (15.632395671810503,24.635197461368666), NE*0.5);
label("$P$", (10.080388928213587,29.181015280559667), NE*0.5);
label("$L$", (7.313369386097323,18.1644843431961), NE*0.5);
label("$K$", (5.175217921734756,27.402217003484928), NE*0.5);
label("$O_1$", (10.763159984060456,24.060232361708145), NE*0.5);
label("$O_2$", (1.3,21.65256600687971), NE*0.5);
\end{asy}
\end{center}
\begin{remark*}
  Questo problema è facilmente bashabile in complessi, consiglio di vedere
  il video di Giuseppe Bassi che spiega come si usano in questo
  problema (\url{https://www.youtube.com/watch?v=MFcWcl5_DDo}) , ho proposto un'altra
  soluzione perchè confesso che con i numeri complessi non ho molta
  dimestichezza.
\end{remark*}

\subsection{G3}


\textbf{Soluzione 1 (Mediana e simmediana)}

Il problema chiede di dimostrare che le rette $AP$ e $AQ$ sono isogonali rispetto
a $\angle BAC$. Poiché $MN\parallel BC$, per l’inverso del teorema di Ceva la retta
$AP$ è la mediana uscente da $A$.
La tesi è quindi equivalente a dimostrare che $AQ$ è la simmediana,
cioè che $AQ$ passa per l’intersezione $T$ delle tangenti alla circonferenza
circoscritta di $ABC$ nei punti $B$ e $C$.

Il punto $Q$ è il punto di Miquel del quadrilatero completo $ANCPBM$,
quindi i quadrilateri $BQNA$ e $CQMA$ sono ciclici. Sia $Y$ la seconda
intersezione di $MN$ con la circonferenza $(BNA)=\omega_b$ e sia $X$ la
seconda intersezione di $MN$ con la circonferenza $(CMA)=\omega_c$.

Si ha $XY\parallel BC$. Inoltre
\[
\dang YXC=\dang MXC=\dang MAC=\dang BAN=\dang BYN=\dang BYX,
\]
quindi $BCYX$ è un trapezio isoscele. Sia $T=BY\cap CX$.
Dal conto degli angoli $T$ è il punto di intersezione delle tangenti in
$B$ e in $C$. Inoltre
\[
\operatorname{Pow}_{\omega_b}(T)=TB\cdot TY=TC\cdot TX=\operatorname{Pow}_{\omega_c}(T),
\]
dunque $T$ appartiene all’asse radicale di $\omega_b$ e $\omega_c$,
che è la retta $AQ$. Quindi $T\in AQ$, come volevasi.

\begin{center}
\begin{asy}
import graph;
size(8cm);

pair A = (6,30);
pair B = (4,22);
pair C = (14,22);
pair M = (5.17753,26.71014);
pair N = (9.28985,26.71014);
pair P = (7.74839,25.33761);
pair Q = (7.39190,22.42187);
pair Y = (1.17391,26.71014);
pair X = (16.82608,26.71014);
pair T = (9,13.66666);
pair K = (9,22);

pen blu = rgb(0.49,0.49,1);
pen viola = rgb(0.29,0,0.51);
pen arancio = rgb(1,0.4,0);
pen rosso = rgb(0.86,0.08,0.24);
pen azzurro = rgb(0.2,0.2,1);

draw(A--B--C--cycle, linewidth(0.4)+blu);
draw(circle((5.45624,24.13820), 2.58699), linewidth(0.8)+viola);
draw(circle((10.77745,23.48759), 3.54933), linewidth(0.8)+viola);
draw(circle((5.23188,25.94202), 4.13002), linewidth(0.8)+arancio);
draw(circle((11.00181,27.00181), 5.83157), linewidth(0.8)+arancio);
draw(X--Y, linewidth(0.8)+blu);
draw(T--A, linewidth(0.8)+rosso);
draw(T--Y, linewidth(0.8)+azzurro);
draw(T--X, linewidth(0.8)+azzurro);
draw(A--K, linewidth(0.8)+rosso);
draw(B--N, linewidth(0.8)+blu);
draw(M--C, linewidth(0.8)+blu);

dot("$A$", A, dir(134));
dot("$B$", B, dir(211));
dot("$C$", C, dir(303));
dot("$M$", M, dir(166));
dot("$N$", N, dir(73));
dot("$P$", P, dir(96));
dot("$Q$", Q, dir(183));
dot("$Y$", Y, dir(166));
dot("$X$", X, dir(21));
dot("$T$", T, dir(43));
dot("$K$", K, dir(43));
\end{asy}
\end{center}



\textbf{Soluzione 2 (Inversione + simmetria rispetto alla bisettrice)}

Come nella Soluzione 1 si ottiene che i quadrilateri $CMQA$ e $BNQA$ sono ciclici.
Si applica un’inversione di centro $A$ e raggio $\sqrt{MA\cdot CA}$,
seguita dalla riflessione rispetto alla bisettrice di $\angle CAB$.
Tale composizione manda la retta $BN$ nella circonferenza $\omega_b$
e la retta $CM$ nella circonferenza $\omega_c$; di conseguenza manda $P$ in $Q$.
Ne segue che le rette $AP$ e $AQ$ sono isogonali.
\section{N}

\subsection{N1}
Noto che $d_n | (100+(n+1)^2)-(100+ n^2)=2n+1$ quindi $d_n| n(2n+1)-2(100+n^2)=
n-200$, e allora $d_n| 2n+1-2(n-200)=401$, quindi $d_n\le 401$. L'uguaglianza la
ottengo scegliendo $n=200$, in quanto $n^2+100 = 100 \cdot 401$, e $(n+1)^2+100 = 101 \cdot
401$.

\subsection{N2}
Sia $D=\{n\in\mathbb{N}:n\mid(2^n+1)\}$.

\textbf{(a)} Sia $p\in D$ primo. Allora $2^p\equiv-1\pmod{p}$. L’ordine $d$ di $2$
modulo $p$ divide $p-1$ (Fermat) e divide $2p$ ma non $p$, dunque $d=2$.
Segue $p\mid3$, quindi $p=3$. Poiché $3\mid9$, l’unico primo in $D$ è $3$.

\textbf{(b)} Sia $p^k\in D$. Lo stesso argomento di (a) dà $p=3$. Per il
lemma di lifting the exponent (applicabile perché l’esponente è dispari)
\[
v_3(2^{3^k}+1)=v_3(2+1)+v_3(3^k)=1+k>k,
\]
quindi $3^k\mid(2^{3^k}+1)$. Tutte le potenze di $3$ stanno in $D$ e
sono le uniche potenze di primi.

\textbf{(d)} Sia $n\in D$ e sia $p$ il più piccolo primo che divide $n$.
Si vede che $n$ è dispari e $2^n\equiv-1\pmod{p}$, dunque l’ordine $d$ di $2$ modulo
$p$ divide $2n$ ma non $n$. Poiché $d\mid(p-1)$ e $p$ è il minimo primo divisor di $n$,
si ha $\gcd(d,n)=1$. Ne segue $d\mid2$, cioè $p=3$. Quindi ognini elemento di $D$ è multiplo di $3$.

\textbf{(c)} Sia $n=pq$ con $p,q$ primi distinti. Per (d) uno dei fattori è $3$;
sia $n=3p$ con $p\neq3$. Allora $2^{3p}\equiv-1\pmod{p}$. Poiché $3p\equiv3\pmod{p-1}$
si ottiene
\[
2^{3p}\equiv2^3\equiv8\pmod{p},
\]
da cui $p\mid9$: contraddizione. Non esistono elementi di $D$ prodotto di due primi distinti.

\textbf{(e)} Sia $n=p^2q$ con $p,q$ primi distinti. Per (d) $3\mid n$.\\

- Se $q=3$ (e $p\neq3$), allora $2^{3p^2}\equiv-1\pmod{p}$. Poiché $3p^2\equiv3\pmod{p-1}$
si ha $2^{3p^2}\equiv8\pmod{p}$, dunque $p\mid9$: contraddizione. \\
- Se $p=3$ (e $q\neq3$), allora $2^{9q}\equiv-1\pmod{q}$.
Analogamente $9q\equiv9\pmod{q-1}$, quindi $2^{9q}\equiv512\equiv-1\pmod{q}$,
ossia $q\mid513=3^3\cdot19$. Dunque $q=19$.

Resta da controllare $171=9\cdot19$. Si ha $v_3(2^{171}+1)=1+v_3(171)=2$,
quindi $9\mid(2^{171}+1)$; inoltre $19\mid(2^{171}+1)$ per la congruenza precedente.
Essendo $9$ e $19$ coprimi, $171\in D$. È l’unico elemento di questa forma.



\end{document}
